\documentclass[11pt]{article}
\usepackage[margin=1.1in]{geometry}
\usepackage{amsmath,amssymb,amsthm}
\usepackage[colorlinks=true,linkcolor=blue,citecolor=blue]{hyperref}

\newtheorem{theorem}{Theorem}[section]
\newtheorem{lemma}[theorem]{Lemma}
\newtheorem{proposition}[theorem]{Proposition}
\newtheorem{corollary}[theorem]{Corollary}
\newtheorem{definition}[theorem]{Definition}
\newtheorem{remark}[theorem]{Remark}

\newcommand{\I}{I}
\newcommand{\Cs}{C_s(\I)}
\newcommand{\Fp}{\mathcal F_p(\I)}
\newcommand{\FDC}{\mathcal F_{DC1}(\I)}
\newcommand{\rsup}{\rho_{\sup}}
\newcommand{\rF}{\rho_{\mathcal F}}
\newcommand{\lm}{\lambda}

\title{Distributional chaos is generic among uniformly convergent\\
non-autonomous systems on the interval}
\author{Yitzchak Shmalo\\
\small Einstein Institute of Mathematics, The Hebrew University of Jerusalem, Jerusalem, Israel\\
\small }
\date{June 2026}

\begin{document}
\maketitle

\begin{abstract}
Balibrea and Ruck\'a [\emph{Genericity of distributional chaos in non-autonomous
systems}, arXiv:2410.10288] proved that in the space of convergent non-autonomous
systems of surjective interval maps, the systems with distributional chaos of type 1
(DC1) form a dense set that is neither open nor closed, and stated explicitly that
the residuality of this set remains an open problem.  We solve the problem in the
canonical complete setting: with the uniform-in-time supremum metric, the DC1
systems contain an \emph{open dense} subset of the space of convergent systems ---
namely, the systems whose limit map carries a strict one-step horseshoe --- and are
therefore residual.  Thus DC1 is a generic property of uniformly convergent
non-autonomous systems on the interval, in contrast with the Cantor set, where the
same authors proved genericity of the absence of distributional chaos.  The proof
combines a quantitative openness--density analysis of strict horseshoes for the limit
map with the inheritance theorem of \v Stef\'ankov\'a and the
Schweizer--Sm\'\i tal characterization of positive entropy on the interval.  Two
complements are included: infinite topological entropy is likewise generic in this
space, by a self-contained common-horseshoe estimate, which removes the
equicontinuity hypothesis from the genericity theorem of Balibrea, Sm\'\i tal and
\v Stef\'ankov\'a (we observe that uniform convergence already implies
equicontinuity of the system); and we delineate precisely what remains open for the
integral-type metrics that distinguish pointwise from uniform convergence, where the
ambient space fails to be complete.
\end{abstract}

\section{Introduction}\label{sec:intro}

A non-autonomous discrete dynamical system (nads) on the compact interval
$\I=[0,1]$ is a sequence $f_{1,\infty}=(f_n)_{n\ge1}$ of continuous surjective maps
$f_n:\I\to\I$, acting through the compositions
$f^n_{1,\infty}=f_n\circ\cdots\circ f_1$.  Following
\cite{BR,BSS}, let $\Cs$ denote the continuous surjections of $\I$ with the
supremum metric $\rsup(f,g)=\sup_x|f(x)-g(x)|$, let
\[
\Fp\;:=\;\Bigl\{(f_n)_{n\ge1}\in\Cs^{\mathbb N}:\ (f_n)\text{ converges in }\rsup
\Bigr\},
\qquad
\rF\bigl(f_{1,\infty},g_{1,\infty}\bigr):=\sup_{n\ge1}\rsup(f_n,g_n),
\]
and write $\lm(f_{1,\infty}):=\lim_nf_n\in C(\I)$ for the limit map.  (As observed
in \cite{BR}, with supremum metrics every convergent system converges uniformly, so
$\Fp$ is the space of \emph{uniformly} convergent systems; we return to other
metrics in Section~\ref{sec:other}.)  A pair $x\ne y$ is DC1-scrambled for
$f_{1,\infty}$ if its upper distributional function is identically $1$ and its lower
distributional function vanishes at some $t_0>0$ (Section~\ref{sec:prelim}); the
system is \emph{DC1} if it has an uncountable set of pairwise DC1-scrambled points.
Let $\FDC\subseteq\Fp$ denote the DC1 systems.

Balibrea and Ruck\'a \cite{BR} proved that $\FDC$ is dense in $\Fp$ (for every
admissible metric), that on the Cantor set the \emph{complement} of distributional
chaos is generic, and that on the interval $\FDC$ is neither open nor closed --- and
they singled out the remaining question (\cite[Section 4]{BR}):

\begin{quote}
``Note that Lemma 2 does not say $\mathcal F_{DC1}(I)$ is not residual.  This
problem stays open.''
\end{quote}

This paper resolves the problem in the canonical complete setting.

\begin{theorem}\label{thm:main}
The set $\FDC$ of DC1 systems contains an open dense subset of
$(\Fp,\rF)$, namely
\[
\lm^{-1}(\mathcal O)
=\bigl\{f_{1,\infty}\in\Fp:\ \lm(f_{1,\infty})\in\mathcal O\bigr\},
\]
where $\mathcal O\subseteq\Cs$ is the (open, dense) set of surjective maps carrying
a strict one-step horseshoe (Definition~\ref{def:horseshoe}).  In particular
$\FDC$ is residual: DC1 is a generic property of uniformly convergent
non-autonomous systems on the interval.
\end{theorem}

The point is that the failure of openness established in \cite[Lemma 2]{BR} is a
boundary phenomenon: a DC1 system may converge to a non-chaotic limit, and such
systems can be tail-perturbed out of $\FDC$.  But systems whose \emph{limit} is
robustly chaotic form an open set --- openness of the horseshoe is measured for the
limit map, where the metric $\rF$ controls $\rsup$-perturbations of the limit ---
and the inheritance theorem of \v Stef\'ankov\'a \cite{Stefankova} pushes the
chaos of the limit back onto every uniformly convergent system above it.  Residual
sets in a complete metric space may well be neither open nor closed; the open dense
\emph{subset} is the correct object, and this observation is the entire content of
the solution.

The same mechanism yields a second genericity statement, this time self-contained.

\begin{theorem}\label{thm:entropy}
The systems with infinite topological entropy contain a dense $G_\delta$ subset of
$(\Fp,\rF)$; indeed for every $m\ge2$ the set of systems whose limit map carries a
strict $m$-branch horseshoe is open and dense, and any such system satisfies
$h(f_{1,\infty})\ge\log m$.  Hence infinite topological entropy is generic in
$\Fp$.
\end{theorem}

Theorem \ref{thm:entropy} recovers the genericity theorem of Balibrea,
Sm\'\i tal and \v Stef\'ankov\'a \cite[Theorem A-type result]{BSS}, which was
stated for \emph{uniformly convergent and equicontinuous} systems, and removes the
second hypothesis --- vacuously, as it turns out:

\begin{proposition}\label{prop:equi}
Every uniformly convergent sequence $(f_n)$ of continuous self-maps of a compact
metric space is an equicontinuous family.  Consequently ``uniformly convergent and
equicontinuous'' defines the same class of systems as ``uniformly convergent''.
\end{proposition}

Finally, Section~\ref{sec:other} discusses the metrics of class $\varrho_1$ of
\cite{BR} (such as the integral metric), which genuinely distinguish pointwise from
uniform convergence.  There the ambient space fails to be complete, the
Baire-category formulation of genericity requires a choice of completion or
framework, and the entropy and chaos transfer between system and limit breaks down
(\cite[Section 3]{BR}); the genericity problem in that setting remains open, and we
state it precisely.

\section{Preliminaries}\label{sec:prelim}

Throughout, $\I=[0,1]$ with the euclidean metric $d$.  For
$f_{1,\infty}=(f_n)\in\Cs^{\mathbb N}$ put $f^0_{1,\infty}=\mathrm{id}$ and
$f^n_{1,\infty}=f_n\circ\cdots\circ f_1$.  For $x,y\in\I$ the lower and upper
distributional functions are
\[
\psi_{xy}(t)=\liminf_{n\to\infty}\frac1n\#\bigl\{0\le j<n:\
d(f^j_{1,\infty}(x),f^j_{1,\infty}(y))<t\bigr\},
\qquad
\psi^*_{xy}(t)=\limsup_{n\to\infty}(\cdots),
\]
and $(x,y)$ is a DC1-scrambled pair if $\psi^*_{xy}\equiv1$ on $(0,\infty)$ and
$\psi_{xy}(t_0)=0$ for some $t_0>0$.  The system is DC1 if some uncountable
$S\subseteq\I$ consists of pairwise DC1-scrambled points.  Topological entropy
$h(f_{1,\infty})$ is defined through $(n,\varepsilon)$-separated sets as in
\cite{KS,BR}; for an autonomous map $f$, $h(f)$ is the usual entropy and we identify
$f$ with the constant system $(f,f,\dots)$.

We use three known results.

\begin{theorem}[{Schweizer--Sm\'\i tal \cite{SS}}]\label{thm:ss}
A continuous interval map has positive topological entropy if and only if it is
distributionally chaotic; in particular $h(f)>0$ implies that $f$ is DC1.
\end{theorem}

\begin{theorem}[{\v Stef\'ankov\'a \cite{Stefankova}}]\label{thm:stef}
Let $(f_n)\subset\Cs$ converge uniformly to $f$.  If $f$ is DC1, then the
non-autonomous system $f_{1,\infty}$ is DC1.
\end{theorem}

\begin{lemma}\label{lem:complete}
$(\Fp,\rF)$ is a complete metric space; hence it is a Baire space, and residuality
in $\Fp$ is meaningful.
\end{lemma}

\begin{proof}
$\Cs$ is complete: a $\rsup$-Cauchy sequence of continuous surjections converges
uniformly to a continuous map, which is surjective because a uniform limit of
surjective self-maps of a compact space is surjective (if $u_k\to u$ and
$u_k(x_k)=y$ with $x_k\to x$ along a subsequence, then $u(x)=y$).  Let
$(f^{(k)}_{1,\infty})_k$ be $\rF$-Cauchy in $\Fp$.  For each $n$ the maps
$f^{(k)}_n$ form a $\rsup$-Cauchy sequence with modulus independent of $n$, hence
converge uniformly to some $g_n\in\Cs$, and $\rF(f^{(k)}_{1,\infty},g_{1,\infty})\to0$.
It remains to check $g_{1,\infty}\in\Fp$, i.e.\ $(g_n)_n$ is $\rsup$-Cauchy.  Given
$\varepsilon>0$ pick $k$ with $\rF(f^{(k)},g)<\varepsilon/3$ and then $N$ with
$\rsup(f^{(k)}_n,f^{(k)}_m)<\varepsilon/3$ for $n,m\ge N$ (possible since
$(f^{(k)}_n)_n$ converges).  Then
$\rsup(g_n,g_m)<\varepsilon$ for $n,m\ge N$.
\end{proof}

\begin{proof}[Proof of Proposition \ref{prop:equi}]
Let $f_n\to f$ uniformly on the compact space $X$ and let $\varepsilon>0$.  Choose
$N$ with $\rsup(f_n,f)<\varepsilon/3$ for $n>N$.  The maps $f,f_1,\dots,f_N$ are
finitely many continuous maps on a compact space, hence uniformly continuous with a
common modulus: there is $\delta>0$ such that $d(x,y)<\delta$ implies
$d(g(x),g(y))<\varepsilon/3$ for $g\in\{f,f_1,\dots,f_N\}$.  For $n>N$ and
$d(x,y)<\delta$,
\[
d(f_n(x),f_n(y))\le d(f_n(x),f(x))+d(f(x),f(y))+d(f(y),f_n(y))<\varepsilon .
\]
Thus the family $(f_n)$ is equicontinuous.
\end{proof}

\section{Strict horseshoes for the limit map}\label{sec:horseshoe}

\begin{definition}\label{def:horseshoe}
Let $m\ge2$.  A map $g\in C(\I)$ carries a \emph{strict $m$-branch horseshoe} with
margin $\theta>0$ if there are pairwise disjoint nonempty closed intervals
$A_1<A_2<\dots<A_m$ in $\I$ such that, with
$[u,v]:=[\min A_1,\max A_m]$ (the convex hull of $\bigcup_iA_i$),
\[
g(A_i)\;\supseteq\;[u-\theta,\,v+\theta]\cap\I
\quad\text{and}\quad
[u-\theta,v+\theta]\subseteq(0,1)
\qquad(1\le i\le m).
\]
Let $\mathcal O_m\subseteq\Cs$ be the set of surjective maps carrying a strict
$m$-branch horseshoe (for some intervals and margin), and $\mathcal O:=\mathcal O_2$.
\end{definition}

\begin{lemma}[Openness]\label{lem:open}
For every $m\ge2$, $\mathcal O_m$ is open in $(\Cs,\rsup)$: if $g$ carries a strict
$m$-horseshoe with margin $\theta$ and $\rsup(h,g)<\theta/2$, then $h$ carries the
same intervals as a strict $m$-horseshoe with margin $\theta/2$.
\end{lemma}

\begin{proof}
$g(A_i)$ is an interval containing $[u-\theta,v+\theta]$, so
$\min_{A_i}g\le u-\theta$ and $\max_{A_i}g\ge v+\theta$.  If
$\rsup(h,g)<\theta/2$ then $\min_{A_i}h\le u-\theta/2$ and
$\max_{A_i}h\ge v+\theta/2$, and $h(A_i)$, being an interval, contains
$[u-\theta/2,v+\theta/2]$.
\end{proof}

\begin{lemma}[Density]\label{lem:dense}
For every $m\ge2$, $\mathcal O_m$ is dense in $(\Cs,\rsup)$.
\end{lemma}

\begin{proof}
Let $g\in\Cs$ and $\varepsilon>0$.  Since $g$ is a continuous surjection of $\I$,
it has a fixed point $z$: the map $x\mapsto g(x)-x$ is $\ge0$ at a preimage of $1$
and $\le0$ at a preimage of $0$.  Replacing $z$ by a nearby point if necessary we
may assume $z\in(0,1)$ (if the only fixed points are $0$ or $1$, say $z=0$, pick
$z'\in(0,1)$ small; then $|g(z')-z'|\le|g(z')-g(0)|+z'\le\omega_g(z')+z'$ is small,
and the construction below around $z'$ works verbatim with this additional error).
We give the construction for a fixed point $z\in(0,1)$; the general case only
changes $\varepsilon$-bookkeeping by the displayed error.

Let $\omega_g$ be the modulus of continuity of $g$, let
$R:=\max\{\eta,\,\omega_g(\eta)\}+\eta$, and choose $\eta>0$ so small
that $J:=[z-\eta,z+\eta]\subseteq(0,1)$ and
\[
\varepsilon':=2R=2\bigl(\max\{\eta,\omega_g(\eta)\}+\eta\bigr)<\varepsilon,
\]
which is possible since $R\to0$ as $\eta\to0$.  Let
$K:=[z-R,\,z+R]\cap\I$; note $g(J)\subseteq[z-\omega_g(\eta),z+\omega_g(\eta)]
\subseteq K$, since $|g(x)-z|=|g(x)-g(z)|\le\omega_g(\eta)$ on $J$.  Define
$h\in C(\I)$ by $h:=g$ off $J$, and on $J$ let $h$ be the piecewise linear map
which fixes the endpoints, $h(z\pm\eta)=g(z\pm\eta)$, and zigzags $2m$ times across
the interval $K$: concretely, partition $J$ into $2m+2$ subintervals; on the first
and last, interpolate linearly from $g(z\mp\eta)$ to the zigzag; on the middle
$2m$ subintervals let $h$ alternate linearly between the values
$\min K$ and $\max K$.

Then $h$ is continuous, and for $x\in J$ both $h(x)$ and $g(x)$ lie in $K$, so
\[
\rsup(h,g)\le\operatorname{diam}K\le2R=\varepsilon'<\varepsilon .
\]
$h$ is surjective: $h(\I)\supseteq g(\I\setminus J)\cup h(J)
=g(\I\setminus J)\cup K\supseteq g(\I\setminus J)\cup g(J)=g(\I)=\I$,
because $g(J)\subseteq K=h(J)$.

Finally $h\in\mathcal O_m$: the $m$ increasing zigzag branches contain pairwise
disjoint closed intervals $A_1<\dots<A_m$ in the interior of $J$ with
$h(A_i)=K$ for each $i$.  Their convex hull $[u,v]$ satisfies
$[u,v]\subseteq J=[z-\eta,z+\eta]$, while $K$ contains
$[z-R,z+R]\cap\I$ and $R-\eta=\max\{\eta,\omega_g(\eta)\}\ge\eta>0$.  Shrink the
intervals $A_i$ slightly so that $[u,v]$ lies in the interior of $J$, and choose
\[
\theta:=\tfrac12\min\bigl\{R-\eta,\ \operatorname{dist}(J,\{0,1\})\bigr\}>0 .
\]
Then $[u-\theta,v+\theta]\subseteq(0,1)$ and
$h(A_i)=K\supseteq[z-R,z+R]\cap\I\supseteq[u-\theta,v+\theta]$, since every point
of $[u-\theta,v+\theta]$ is within $\eta+\theta\le R$ of $z$ and lies in $(0,1)$.
Thus $h$ carries a strict $m$-branch horseshoe with margin $\theta$.
\end{proof}

\begin{remark}
The proof shows more: the horseshoe may be implanted at any prescribed scale
$\eta$, and the margin obtained is at least $\eta$.  For $m=2$ this is the
classical fact that ``robust horseshoes are $C^0$-dense''; we could not locate a
reference in the surjective category with explicit margins, so we included the
proof.  Compare Kawaguchi \cite{Kawaguchi}, who proved $C^0$-density of DC1 maps.
\end{remark}

\section{Proofs of the main theorems}\label{sec:proofs}

\begin{lemma}\label{lem:lambda}
The limit-map function $\lm:(\Fp,\rF)\to(\Cs,\rsup)$ is $1$-Lipschitz.
\end{lemma}

\begin{proof}
$\rsup\bigl(\lm(f_{1,\infty}),\lm(g_{1,\infty})\bigr)
=\lim_n\rsup(f_n,g_n)\le\rF(f_{1,\infty},g_{1,\infty})$.
\end{proof}

\begin{lemma}\label{lem:preimage-dense}
For every $m\ge2$, the set $\lm^{-1}(\mathcal O_m)$ is open and dense in
$(\Fp,\rF)$.
\end{lemma}

\begin{proof}
Openness is immediate from Lemmas \ref{lem:open} and \ref{lem:lambda}.  For
density, let $f_{1,\infty}\in\Fp$ with limit $f$ and let $\varepsilon>0$.  By
Lemma~\ref{lem:dense} pick $g\in\mathcal O_m$ with $\rsup(g,f)<\varepsilon/2$, and
pick $N$ with $\rsup(f_n,f)<\varepsilon/2$ for $n>N$.  Define
$g_{1,\infty}$ by $g_n:=f_n$ for $n\le N$ and $g_n:=g$ for $n>N$.  Then
$g_{1,\infty}\in\Fp$ (its tail is constant), $\lm(g_{1,\infty})=g\in\mathcal O_m$,
and
\[
\rF(f_{1,\infty},g_{1,\infty})
=\sup_{n>N}\rsup(f_n,g)
\le\sup_{n>N}\bigl[\rsup(f_n,f)+\rsup(f,g)\bigr]<\varepsilon .
\]
\end{proof}

\begin{lemma}[Common horseshoe forces entropy]\label{lem:common}
Let $f_{1,\infty}\in\Cs^{\mathbb N}$ and suppose there are $m\ge2$, $N\ge0$,
pairwise disjoint closed intervals $A_1<\dots<A_m$ with convex hull $[u,v]$, and
$\theta>0$ such that
\[
f_n(A_i)\supseteq[u,v]
\qquad\text{for all }n>N,\ 1\le i\le m .
\]
Then $h(f_{1,\infty})\ge\log m$.
\end{lemma}

\begin{proof}
Let $\gamma>0$ be the minimal gap between consecutive intervals $A_i$.  Fix a
horizon $L\ge1$ and a word $\omega\in\{1,\dots,m\}^L$.  Choose
$y_L\in A_{\omega_L}$ arbitrarily.  Descending from $j=L-1$ to $j=1$: given
$y_{j+1}\in A_{\omega_{j+1}}\subseteq[u,v]$, the covering hypothesis (applied to
$n=N+j>N$) gives $f_{N+j}(A_{\omega_j})\supseteq[u,v]\ni y_{j+1}$, and since the
continuous image of an interval is an interval there is
$y_j\in A_{\omega_j}$ with $f_{N+j}(y_j)=y_{j+1}$.  This produces
$y_1\in A_{\omega_1}$ whose forward orbit under the tail
$(f_{N+1},f_{N+2},\dots)$ satisfies
$f_{N+j-1}\circ\cdots\circ f_{N+1}(y_1)=y_j\in A_{\omega_j}$ for $j=1,\dots,L$.
Since each $f_n$ is surjective, so is $f^N_{1,\infty}$; choose
$x_\omega$ with $f^N_{1,\infty}(x_\omega)=y_1$.  Then
$f^{N+j-1}_{1,\infty}(x_\omega)=y_j\in A_{\omega_j}$ for $j=1,\dots,L$.

If $\omega\ne\omega'$ differ at position $j$, the orbits of $x_\omega,x_{\omega'}$
at time $N+j-1$ lie in distinct intervals $A_{\omega_j}\ne A_{\omega'_j}$, hence
are $\gamma$-separated.  Thus $\{x_\omega\}$ is an
$(N+L,\gamma)$-separated set of cardinality $m^L$, and
\[
h(f_{1,\infty})\ \ge\ \limsup_{L\to\infty}\frac{\log m^L}{N+L}\ =\ \log m .
\]
\end{proof}

\begin{proof}[Proof of Theorem \ref{thm:main}]
By Lemma~\ref{lem:preimage-dense} (with $m=2$), $\lm^{-1}(\mathcal O)$ is open and
dense in $\Fp$.  Let $f_{1,\infty}\in\lm^{-1}(\mathcal O)$ with limit $f$.  By
Lemma~\ref{lem:common} applied to the constant system $(f,f,\dots)$ --- the
covering hypothesis holds with $N=0$ since $f(A_i)\supseteq[u-\theta,v+\theta]
\supseteq[u,v]$ --- we get $h(f)\ge\log2>0$.  By Theorem~\ref{thm:ss}, $f$ is DC1.
By Theorem~\ref{thm:stef}, the uniformly convergent system $f_{1,\infty}$ is DC1.
Hence $\lm^{-1}(\mathcal O)\subseteq\FDC$, and $\FDC$, containing an open dense
set, is residual in the complete space $(\Fp,\rF)$ (Lemma~\ref{lem:complete}).
\end{proof}

\begin{proof}[Proof of Theorem \ref{thm:entropy}]
Fix $m\ge2$ and let $f_{1,\infty}\in\lm^{-1}(\mathcal O_m)$, with limit $f$
carrying a strict $m$-horseshoe on intervals $A_1<\dots<A_m$ with margin $\theta$.
Choose $N$ with $\rsup(f_n,f)<\theta/2$ for $n>N$.  By the proof of
Lemma~\ref{lem:open}, every $f_n$ with $n>N$ satisfies
$f_n(A_i)\supseteq[u-\theta/2,v+\theta/2]\supseteq[u,v]$.  By
Lemma~\ref{lem:common}, $h(f_{1,\infty})\ge\log m$.  Hence
\[
\bigcap_{m\ge2}\lm^{-1}(\mathcal O_m)\subseteq
\bigl\{f_{1,\infty}:\ h(f_{1,\infty})=\infty\bigr\},
\]
and the left-hand side is a countable intersection of open dense sets
(Lemma~\ref{lem:preimage-dense}), hence a dense $G_\delta$.
\end{proof}

\begin{remark}\label{rem:bss}
Theorem~\ref{thm:entropy} together with Proposition~\ref{prop:equi} subsumes the
genericity of infinite entropy for ``uniformly convergent and equicontinuous''
systems from \cite{BSS}: by Proposition~\ref{prop:equi} the equicontinuity
hypothesis is implied by uniform convergence, and the proof above is independent
of \cite{BSS}.  Note also that all three results
(Theorems \ref{thm:main}, \ref{thm:entropy} and the density theorem of \cite{BR})
now align with the autonomous interval theory: positive entropy, DC1, and their
genericity transfer along the limit map.
\end{remark}

\begin{remark}[Contrast with the Cantor set]
On the Cantor set, Balibrea and Ruck\'a proved that the \emph{absence} of
distributional chaos is generic for convergent systems \cite[Lemma 1]{BR},
following Darji--D'Aniello \cite{DD}.  Theorem~\ref{thm:main} shows the interval
behaves in the opposite way, completing the table of \cite[Section 4]{BR}: DC1 is
generic in $\Fp(\I)$ and non-DC1 is generic in $\Fp(Q)$.  The mechanism of the
difference is visible in the proofs: on the interval, surjectivity plus the
intermediate value theorem make strict horseshoes of the limit map both robust and
densely implantable, while on the Cantor set the clopen structure makes
adding-machine-like (chaos-free) limits robust and dense.
\end{remark}

\section{The metrics distinguishing pointwise from uniform convergence}
\label{sec:other}

Balibrea and Ruck\'a formulate their results for arbitrary metrics on $\Cs$,
split into the class $\varrho_2$ (metrics, like $\rsup$, for which convergence of
maps forces uniform convergence) and the class $\varrho_1$ (metrics, like the
integral metric $\rho_{\mathrm{int}}(f,g)=\int_0^1d(f(x),g(x))\,dx$, for which
genuinely pointwise-but-not-uniformly convergent sequences of maps exist)
\cite[Remark 1]{BR}.  Theorems \ref{thm:main} and \ref{thm:entropy} settle the
genericity problem for the supremum metric, the canonical representative of
$\varrho_2$.  For other metrics in $\varrho_2$ the underlying set of systems is
the same --- every $\varrho_2$-convergent system is uniformly convergent --- but
residuality depends on the topology of the system space, so each such metric for
which the space is Baire requires checking that the horseshoe class
$\lm^{-1}(\mathcal O)$ remains open and dense; for metrics uniformly equivalent
to $\rsup$ this is automatic.

For $\varrho_1$-metrics the situation is genuinely different, for three reasons
already implicit in \cite{BR}.  First, $(\Cs,\rho_{\mathrm{int}})$ is not complete
--- pointwise limits of continuous surjections need not be continuous --- so the
system space is not complete either, and \cite{BR} must explicitly demand
continuity of limit maps; the Baire category theorem is then unavailable without
either passing to a completion (changing the space) or verifying the Baire
property by hand.  Second, the entropy bridge fails: there are systems with
positive (even infinite) entropy converging pointwise to the identity
\cite[Figure 1]{BR}, so the limit map no longer certifies the dynamics of the
system, and our key mechanism --- openness of chaos \emph{at the level of the
limit} --- has no analogue.  Third, the inheritance theorem
(Theorem~\ref{thm:stef}) is a theorem about uniform convergence.  We therefore
record what remains of the original question:

\begin{quote}
\textbf{Open problem.}  Let $\rho\in\varrho_1$ (for instance
$\rho=\rho_{\mathrm{int}}$), and consider the set $F$ of pointwise convergent
systems of surjective interval maps with continuous limit, with the metric
$\sup_n\rho(f_n,g_n)$.  Is the subspace $F$ Baire?  If so, is DC1 generic in it?
\end{quote}

Our results show that any positive answer must use a mechanism that does not
factor through the limit map; the examples of \cite[Section 3]{BR} (positive
entropy, eventually-fixed-point systems converging pointwise to the identity)
suggest that smallness at the level of the limit and largeness at the level of the
system can coexist on residual-looking sets, and we expect the answer to depend on
the metric within $\varrho_1$.

\subsection*{Acknowledgments}
The research presented in this paper was supported by the European Research Council
(ERC) under the European Union's Horizon Europe research and innovation programme
(grant agreement No.\ 101041711), by the Israel Science Foundation (grant number
2258/19) and by the Israel Science Foundation (ISF Grant 4101/25).

\subsection*{Declaration on generative AI}
During the preparation of this work the author used generative AI tools to
accelerate drafting and revision.  All mathematical claims, proofs, and
bibliographic information were subsequently reviewed and edited by the author, who
takes full responsibility for the content of the manuscript.

\begin{thebibliography}{9}

\bibitem{BR}
F.~Balibrea and L.~Ruck\'a,
\emph{Genericity of distributional chaos in non-autonomous systems},
preprint, arXiv:2410.10288 (2024); see also
\emph{Density of distributional chaos in non-autonomous systems},
Topology Appl.\ \textbf{380} (2026), 109735.

\bibitem{BSS}
F.~Balibrea, J.~Sm\'\i tal and M.~\v Stef\'ankov\'a,
\emph{On generic properties of nonautonomous dynamical systems},
Internat.\ J.\ Bifur.\ Chaos Appl.\ Sci.\ Engrg.\ \textbf{28} (2018), no.~8,
1850102.

\bibitem{DD}
E.~D'Aniello and U.~B.~Darji,
\emph{Chaos among self-maps of the Cantor space},
J.\ Math.\ Anal.\ Appl.\ \textbf{381} (2011), 781--788.

\bibitem{Kawaguchi}
N.~Kawaguchi,
\emph{Distributionally chaotic maps are $C^0$-dense},
Proc.\ Amer.\ Math.\ Soc.\ \textbf{147} (2019), no.~12, 5339--5348.

\bibitem{KS}
S.~Kolyada and L'.~Snoha,
\emph{Topological entropy of nonautonomous dynamical systems},
Random Comput.\ Dynam.\ \textbf{4} (1996), 205--233.

\bibitem{SS}
B.~Schweizer and J.~Sm\'\i tal,
\emph{Measures of chaos and a spectral decomposition of dynamical systems on the
interval},
Trans.\ Amer.\ Math.\ Soc.\ \textbf{344} (1994), no.~2, 737--754.

\bibitem{Stefankova}
M.~\v Stef\'ankov\'a,
\emph{Inheriting of chaos in uniformly convergent nonautonomous dynamical systems
on the interval},
Discrete Contin.\ Dyn.\ Syst.\ \textbf{36} (2016), no.~6, 3435--3443.

\end{thebibliography}

\end{document}
